\section{The operator-product estimates on the selected history}
\label{sec:ope-history-service}

The local projection and the several-insertion expansion must use the
same adjacent transition. We establish this compatibility on the family
of Section~\ref{sec:selected-centered-service}. Throughout this section
the profiles have support exponent $100$, as selected in
Section~\ref{sec:ope-scale-purchases}. The orders and spatial powers
still increase. No background representation is chosen again.

\subsection{Finite elementary representations of transferred fields}

Fix a finite torus, a representation $R$ and an anchor $o$. For a support
$X$ and a nonnegative bound $h$, write
\[
                  m_R(X,h)=e^{b C_B(X)}h,
\]
where $b\ge0$ and $B$ are its support exponent and covering radius.
The reserved regular slot on each nonempty support and the reserved
forced slot at each admissible support--inventory pair are retained
from the marked construction. A finite elementary representation of
an ordinary part $g$ and a forced profile $K$ means identities
\begin{equation}\label{eq:ope-finite-elementary}
 g=p+\sum_{j\in I}z_j\operatorname{ord}(M_j),\qquad
 K=\sum_{j\in I}z_j\operatorname{flag}(M_j),\qquad
 \sum_{j\in I}|z_j|\operatorname{cost}(M_j)\le A.
\end{equation}
Here $I$ is finite, all insertion supports contain $o$, and $p\in\mathbb C$
is the separately recorded constant. In particular $K$ is the actual
profile, not merely an array with the same integrated density.

Suppose first that $g=\sum_i f_i$, where the sum is countable, each
$f_i$ is measurable and gauge invariant, depends only on its support
$X_i$ and the flags there, vanishes unless $X_i$ is healthy, and
$o\in X_i$. Assume
\[
       |f_i(U)|\le h_i,\qquad
                   \sum_i m_R(X_i,h_i)\le A.
\]
Since $b\ge0$, $\sum_i h_i<\infty$. Group by the finitely many
supports of the torus and put
\[
 f_X=\sum_{i:X_i=X}f_i,\qquad h_X=\sum_{i:X_i=X}h_i.
\]
These series converge uniformly and absolutely. Their sums preserve
measurability, gauge invariance, local dependence and the healthy-set
condition. The real and imaginary parts of $f_X$ can be placed in
the reserved regular slot of $X$, each with bound $h_X$. Their
coefficients are $1$ and $i$. The resulting elementary representation
has ordinary part $f_X$, zero forced profile and price at most
$2m_R(X,h_X)$. Zero functions may be omitted. Summing over the finite
support set gives \eqref{eq:ope-finite-elementary} with $p=0$, $K=0$
and price $2A$. The equality of the price sums follows by grouping
an absolutely convergent nonnegative series; its support weight is
constant within each group.

The same argument applies to a dominated forced array. Group its
coefficients by their support--inventory pair $(X,J)$. Suppose the
grouped coefficient vanishes unless the flags in $X$ are exactly $J$,
where $J$ is nonempty, and every nonzero group has $o\in X$.
Use the reserved forced slot at $(X,J)$ and split the grouped function
into real and imaginary parts. An array with total weighted majorant
$A$ therefore has a finite elementary representation of its profile,
with ordinary part zero and price at most $2A$. The exact inventory
condition is preserved under grouping. This construction uses only
finitely many pairs even when the original array has infinitely many
serial addresses.

We apply this compression to an elementary insertion $M$ at a fine
scale, with mass $m=\operatorname{cost}(M)\ge0$. Denote its block
transfer by $\mathcal T$ and the forced derivative fixed in
Section~\ref{sec:projection-payment-producer} by $\mathcal D M$.
The exact differentiated expansion gives an ordinary sum and this
forced derivative. Its nonprincipal ordinary mass is at most $p_M$;
its possible principal term has mass at most $E_*m$; and its forced
mass is at most $q_M$. The local derivative and support estimates give
\begin{equation}\label{eq:ope-elementary-compression-bills}
           2p_M\le m,\qquad 2q_M\le m,\qquad E_*\le L.
\end{equation}
Here $E_*$ is the fixed exponential charge for enlarging the inherited
support. The last inequality is the geometric block-length floor in
Section~\ref{sec:projection-payment-producer}. More precisely, $p_M$
is the nonprincipal regular bill of that construction, and $q_M$
is its forced derivative bill before division by the clipping gain.
Put $\kappa=\sqrt{\text{coarse forced allowance}}$ and let $H_0$
be the degree-zero splitting factor. Its explicit form is
\[
 H_0=1+2|v_0|(1+C_{\rm orb}+|f_0|)\ge1,
 \qquad c_0=L^{-2/5}/H_0\le1,
\]
where $v_0,f_0$ are the fixed finite local-splitting constants and
$C_{\rm orb}\ge0$ is its motion price. The two degree-zero payments
are exactly
\[
          2p_M\le c_0m,\qquad
          2\kappa^{-1}q_M\le c_0m,
                       \qquad 0<\kappa\le1.
\]
They imply the first two inequalities in
\eqref{eq:ope-elementary-compression-bills}. This uses the same
unscaled forced derivative as the profile identity. The principal contribution
is absent for a forced input. For an ordinary input its masked value
is at most its original height: the single-insertion, order-zero
cluster has no additional insertion factor. Enlarging its support
therefore costs precisely $E_*$ times its original mass. All remaining
clusters are charged to $p_M$.

The compression just proved gives a finite elementary representation
of the transferred ordinary part and the profile of $\mathcal D M$
with price at most
\[
 2(p_M+E_*m)+2q_M\le(2+2L)m\le6Lm,
\]
where $L\ge1$. Its inserted density is exactly $\mathcal T$ applied
to the density of $M$, by the differentiated block identity. Thus the
profile as well as the density is retained.

This conclusion extends to every finite representation
\eqref{eq:ope-finite-elementary}. Apply it to each $M_j$ and multiply
its output coefficients by $z_j$. Concatenating the finite lists
gives total price at most $6LA$. The constant term remains $p$
because $\mathcal T\rho_f=\rho_c$. Finite sums commute with integration:
each elementary density is integrable under the lawful representation,
and the transition is the already constructed kernel. The output
profile is exactly $\sum_jz_j\mathcal P(\mathcal D M_j)$. This proves
the transport bound needed when a grouped observable is carried
through intermediate scales.

\subsection{Several insertions and their common allowances}

At adjacent levels $m+1,m$ set
\[
 x=\widehat\beta_m,\quad y=\widehat\beta_{m+1},\quad
 L=L_{J(m+1)},\quad N=N_{J(m+1)},\quad
 \delta_m=y^{-2},\quad \sigma_m=y^{-1}.
\]
Let $\chi_m$ be the forced-gas price from the two successive
absorptions, with its actual positive suppression constant and this
block length. It is the price denoted by $t$ in
Section~\ref{sec:finite-menu-history-handoff}. The single-insertion
and arbitrary finite-menu estimates of that section use this same
$\mathcal D$, with input support exponent $96$, ordinary output
exponent $100+1/2$, and forced output exponent $101$. Their output
allowances are respectively $1$ and $\chi_m$.

To apply them to several slots in each replica cell, take the finite
marker set to be the product of the cell set and its slot set.
The compiler of the finite-menu construction already permits every
finite marker set. The ordinary coefficient is indexed by individual
markers and the forced coefficients by their subsets. Grouping by
the cell coordinate is exactly the substitution used in
Section~\ref{sec:ope-grouping-algebra}. The quiet coefficient stays
the actual background profile. Before grouping, the singleton coefficient
for each individual marker (cell, slot) is the sum of the derivatives
of its original insertions through $\mathcal D$. After grouping,
the coefficient of a singleton cell is the old coefficient at that
cell's full slot set, by the definition of the regrouping map in
Section~\ref{sec:ope-grouping-algebra}. For example, with two slots,
the coefficient of $e_{\{1,2\}}$ in
$\exp(z_1e_{\{1\}}+z_2e_{\{2\}})$ is $z_1z_2$;
this is the resulting singleton-cell coefficient. The motion identities
follow by applying the same
permutation to the cell coordinate and leaving the slot coordinate
fixed; equal weights on corresponding slots preserve every weighted
estimate. Thus increasing the number of slots does not change the
background transition or its insertion transport.

Write $B=B_m$ and $\alpha=100+1/2$. The enlarged-support bound of
Section~\ref{sec:ope-grouping} gives the local price
\begin{equation}\label{eq:ope-history-cluster-price}
 s_m=2(1+N)\delta_m(Ne^\alpha+1)+(1+N)\chi_m.
\end{equation}
The two copies in its first term are the healthy and unhealthy
pieces of the ordinary grouped entry. The last term is the forced
entry. The purchases in
\eqref{eq:projection-occupancy-payments} give, simultaneously,
\begin{gather}
 0<\delta_m\le1,\qquad \chi_m\le1,\qquad
 (2B_m+1)^4s_m\le\tfrac12,\qquad e^{s_m}-1\le1,
 \label{eq:ope-history-smallness}\\
 \alpha(2048B_m+1)^4\le\sqrt{\widehat\beta_{m+1}}.
 \label{eq:ope-history-support-cost}
\end{gather}
The allowance for the elementary degree-$d$ projection is
$X_m=\exp(\sqrt{\widehat\beta_{m+1}})\ge1$, by
\eqref{eq:projection-large-allowance}. Both that projection and the
$6L$ elementary transport use the same $\mathcal D$ just identified.

\subsection{Uniformity above the reference level}

For each grade $d$ retain the reference level $t_d$ with
$d+3\le J(t_d)$. If $m\ge t_d$, monotonicity gives
$d\le J(m)$. The numerical purchases were made for every degree
up to the predecessor stage, and for every permitted adjacent pair
of stages, before the scale sequence was realized. Hence the
degree-$d$ projection is available at every such $m$, for every
cutoff $N_0>m$ and every period. The constants in
\eqref{eq:ope-history-smallness}--\eqref{eq:ope-history-support-cost}
depend only on $m$, not on $N_0$ or the period.

The adjacent transition retained by the selected history is the one
at the chosen point of its interval family. Substituting the equalities
identifying its fine and coarse representations with the corresponding
history entries changes none of the formulas above. Consequently the
ordinary menu estimates, the several-slot estimates, the finite
elementary transport, and the final local projection hold together
on that history. These supply the local hypotheses used in
Sections~\ref{sec:ope-local-projection} and
\ref{sec:ope-grouping}. The later continuum limiting and cutoff
arguments use their previously established estimates; they require
no new choice of background history.
